\documentclass[12pt, titlepage]{article}

\usepackage{amsmath, mathtools}

\usepackage[round]{natbib}
\usepackage{amsfonts}
\usepackage{amssymb}
\usepackage{graphicx}
\usepackage{colortbl}
\usepackage{xr}
\usepackage{hyperref}
\usepackage{longtable}
\usepackage{xfrac}
\usepackage{tabularx}
\usepackage{float}
\usepackage{siunitx}
\usepackage{booktabs}
\usepackage{multirow}
\usepackage[section]{placeins}
\usepackage{caption}
\usepackage{fullpage}

\hypersetup{
bookmarks=true,     % show bookmarks bar?
colorlinks=true,       % false: boxed links; true: colored links
linkcolor=red,          % color of internal links (change box color with linkbordercolor)
citecolor=blue,      % color of links to bibliography
filecolor=magenta,  % color of file links
urlcolor=cyan          % color of external links
}

\usepackage{array}

\externaldocument{../../SRS/SRS}
\externaldocument{../SoftArchitecture/MG}

\input{../../Comments.text}
\input{../../Common.text}

\newcommand{\mref}[1]{M\ref{#1}}

\begin{document}

\title{Module Interface Specification for \progname{}}

\author{\authname}

\date{\today}

\maketitle

\pagenumbering{roman}

\section{Revision History}

\begin{tabularx}{\textwidth}{p{3cm}p{2cm}X}
\toprule {\bf Date} & {\bf Version} & {\bf Notes}\\
\midrule
Date 1 & 1.0 & Notes\\
Date 2 & 1.1 & Notes\\
\bottomrule
\end{tabularx}

~\newpage

\section{Symbols, Abbreviations and Acronyms}

See SRS Documentation at \wss{give url}

\wss{Also add any additional symbols, abbreviations or acronyms}

\newpage

\tableofcontents

\newpage

\pagenumbering{arabic}

\section{Introduction}

The following document details the Module Interface Specifications for
\wss{Fill in your project name and description}

Complementary documents include the System Requirement Specifications
and Module Guide.  The full documentation and implementation can be
found at \url{...}.  \wss{provide the url for your repo}

\section{Notation}

\wss{You should describe your notation.  You can use what is below as
  a starting point.}

The structure of the MIS for modules comes from \citet{HoffmanAndStrooper1995},
with the addition that template modules have been adapted from
\cite{GhezziEtAl2003}.  The mathematical notation comes from Chapter 3 of
\citet{HoffmanAndStrooper1995}.  For instance, the symbol := is used for a
multiple assignment statement and conditional rules follow the form $(c_1
\Rightarrow r_1 | c_2 \Rightarrow r_2 | ... | c_n \Rightarrow r_n )$.

The following table summarizes the primitive data types used by \progname. 

\begin{center}
\renewcommand{\arraystretch}{1.2}
\noindent 
\begin{tabular}{l l p{7.5cm}} 
\toprule 
\textbf{Data Type} & \textbf{Notation} & \textbf{Description}\\ 
\midrule
character & char & a single symbol or digit\\
integer & $\mathbb{Z}$ & a number without a fractional component in (-$\infty$, $\infty$) \\
natural number & $\mathbb{N}$ & a number without a fractional component in [1, $\infty$) \\
real & $\mathbb{R}$ & any number in (-$\infty$, $\infty$)\\
\bottomrule
\end{tabular} 
\end{center}

\noindent
The specification of \progname \ uses some derived data types: sequences, strings, and
tuples. Sequences are lists filled with elements of the same data type. Strings
are sequences of characters. Tuples contain a list of values, potentially of
different types. In addition, \progname \ uses functions, which
are defined by the data types of their inputs and outputs. Local functions are
described by giving their type signature followed by their specification.

\section{Module Decomposition}

The following table is taken directly from the Module Guide document for this project.

\begin{table}[h!]
\centering
\begin{tabular}{p{0.3\textwidth} p{0.6\textwidth}}
\toprule
\textbf{Level 1} & \textbf{Level 2}\\
\midrule

{Hardware-Hiding} & ~ \\
\midrule

\multirow{7}{0.3\textwidth}{Behaviour-Hiding} & Input Parameters\\
& Output Format\\
& Output Verification\\
& Temperature ODEs\\
& Energy Equations\\ 
& Control Module\\
& Specification Parameters Module\\
\midrule

\multirow{3}{0.3\textwidth}{Software Decision} & {Sequence Data Structure}\\
& ODE Solver\\
& Plotting\\
\bottomrule

\end{tabular}
\caption{Module Hierarchy}
\label{TblMH}
\end{table}

\newpage
~\newpage

\section{MIS of \wss{Module Name}} \label{Module} \wss{Use labels for
  cross-referencing}

\wss{You can reference SRS labels, such as R\ref{R_Inputs}.}

\wss{It is also possible to use \LaTeX for hypperlinks to external documents.}

\subsection{Module}

\wss{Short name for the module}

\subsection{Uses}


\subsection{Syntax}

\subsubsection{Exported Constants}

\subsubsection{Exported Access Programs}

\begin{center}
\begin{tabular}{p{2cm} p{4cm} p{4cm} p{2cm}}
\hline
\textbf{Name} & \textbf{In} & \textbf{Out} & \textbf{Exceptions} \\
\hline
\wss{accessProg} & - & - & - \\
\hline
\end{tabular}
\end{center}

\subsection{Semantics}

\subsubsection{State Variables}

\wss{Not all modules will have state variables.  State variables give the module
  a memory.}

\subsubsection{Environment Variables}

\wss{This section is not necessary for all modules.  Its purpose is to capture
  when the module has external interaction with the environment, such as for a
  device driver, screen interface, keyboard, file, etc.}

\subsubsection{Assumptions}

\wss{Try to minimize assumptions and anticipate programmer errors via
  exceptions, but for practical purposes assumptions are sometimes appropriate.}

\subsubsection{Access Routine Semantics}

\noindent \wss{accessProg}():
\begin{itemize}
\item transition: \wss{if appropriate} 
\item output: \wss{if appropriate} 
\item exception: \wss{if appropriate} 
\end{itemize}

\wss{A module without environment variables or state variables is unlikely to
  have a state transition.  In this case a state transition can only occur if
  the module is changing the state of another module.}

\wss{Modules rarely have both a transition and an output.  In most cases you
  will have one or the other.}

\subsubsection{Local Functions}

\wss{As appropriate} \wss{These functions are for the purpose of specification.
  They are not necessarily something that is going to be implemented
  explicitly.  Even if they are implemented, they are not exported; they only
  have local scope.}

% =====================================================================
\section{MIS of Request Lifecycle (\mref{mLifecycle})} \label{misLifecycle}
% =====================================================================

\subsection{Module}

Request Lifecycle (\mref{mLifecycle}). An abstract data type whose instances are
inference requests, each identified by a \texttt{RequestId}. It is implemented
in the on-chain program.

\subsection{Uses}

\begin{itemize}
\item \textbf{Provider Registry (\mref{mRegistry}):}
  \texttt{getTaskParameters}, \texttt{isRegistered}, \texttt{getReputation},
  \texttt{setReputation}, \texttt{getStake}, \texttt{getReservedStake},
  \texttt{reserveStake}, \texttt{releaseStake}, \texttt{deductStake},
  \texttt{removeProvider}, \texttt{isValidOutput}, \texttt{getSlashHistory},
  \texttt{recordParticipation}, and the constant
  \texttt{MINIMUM\_PROVIDER\_STAKE}. (\texttt{TaskParameters} is a record
  with \texttt{minProviders}, \texttt{maxProviders}, and
  \texttt{acceptancePeriod}.)
\item \textbf{Prediction Aggregator (\mref{mAggregator}):}
  \texttt{weightedMedian}, \texttt{resultSpread}, and \texttt{AggregationError}.
\item \textbf{Reputation Handler (\mref{mReputation}):}
  \texttt{updateReputations} and \texttt{forfeitStake}.
\item \textbf{Rewards Calculator (\mref{mRewards}):}
  \texttt{escrowPayment}, \texttt{distributeRewards}, \texttt{retryReward},
  \texttt{forfeitReward}, \texttt{refundPayment}, \texttt{payCompensation},
  and \texttt{TransferFailedError}.
\item \textbf{Penalty Enforcer (\mref{mPenalty}):}
  \texttt{applyPenalty}, \texttt{collectCompensation}, \texttt{PenaltyDecision},
  and \texttt{ParticipationOutcome}.
\end{itemize}

This module calls no Token Manager (\mref{mToken}) routine. Every movement of
tokens for a request happens inside Rewards Calculator, Penalty Enforcer, or
Reputation Handler.

The routine names above are the interface this MIS assumes from the other
modules; they must match those modules' MIS.

\subsection{Syntax}

\subsubsection{Exported Constants}
\begin{itemize}
\item \texttt{COMMIT\_PERIOD}: Duration -- maximum length of the commit phase,
  measured from the end of the acceptance period
\item \texttt{REVEAL\_PERIOD}: Duration -- maximum length of the reveal phase,
  measured from the end of the commit phase
\item \texttt{PENDING\_TRANSFER\_PERIOD}: Duration -- time after the end of the
  reveal phase during which unresolved reward transfers may be retried
\item \texttt{MINIMUM\_REQUEST\_PAYMENT}: $\mathbb{N}$ -- the lowest payment a
  request may offer (payments must exceed it)
\end{itemize}

\subsubsection{Exported Types}
\begin{itemize}
\item \texttt{RequestStatus} = \{Acceptance, Commit, Reveal, Finalization,
  PendingTransfer, Completed, Failed, Cancelled\}
\item \texttt{RevealResult} = \{Accepted, Mismatch, NonConforming\}
\item \texttt{Hash}: fixed-length byte string; \texttt{Salt}: byte string
\end{itemize}

\subsubsection{Exported Access Programs}

\begin{longtable}{p{3.3cm} p{4.2cm} p{3.1cm} p{5cm}}
\hline
\textbf{Name} & \textbf{In} & \textbf{Out} & \textbf{Exceptions} \\
\hline
\endfirsthead
\hline
\textbf{Name} & \textbf{In} & \textbf{Out} & \textbf{Exceptions} \\
\hline
\endhead
computeCommitment & RequestId, Address, $\mathbb{R}$, Salt & Hash & - \\
\hline
createRequest & TaskId, Address, $\mathbb{N}$ & RequestId &
  UnknownTaskError, InvalidPaymentError, TransferFailedError \\
\hline
acceptRequest & RequestId, Address & - &
  UnknownRequestError, WrongPhaseError, DeadlinePassedError,
  NotRegisteredError, AlreadyAcceptedError, InsufficientStakeError \\
\hline
submitCommitment & RequestId, Address, Hash & - &
  UnknownRequestError, WrongPhaseError, DeadlinePassedError,
  NotAcceptedError, AlreadyCommittedError \\
\hline
revealPrediction & RequestId, Address, $\mathbb{R}$, Salt & RevealResult &
  UnknownRequestError, WrongPhaseError, DeadlinePassedError,
  NoCommitmentError, AlreadyRevealedError \\
\hline
advancePhase & RequestId & RequestStatus &
  UnknownRequestError, ArithmeticOverflowError \\
\hline
retryTransfer & RequestId, Address & - &
  UnknownRequestError, WrongPhaseError, DeadlinePassedError,
  NoUnresolvedTransferError, TransferFailedError \\
\hline
getStatus & RequestId & RequestStatus & UnknownRequestError \\
\hline
getResult & RequestId & $\mathbb{R} \times \mathbb{R} \times \mathbb{R}
  \times \mathbb{N}$ (result, minimum, maximum, count) &
  UnknownRequestError, ResultUnavailableError \\
\hline
getDeadlines & RequestId & Timestamp $\times$ Timestamp $\times$ Timestamp
  $\times$ Timestamp & UnknownRequestError \\
\hline
getUnresolvedTransfers & RequestId & set of (Address $\times$ $\mathbb{N}$) &
  UnknownRequestError \\
\hline
\end{longtable}

\subsection{Semantics}

\subsubsection{State Variables}

Each request \texttt{r} has the following state variables.

\begin{itemize}
\item \texttt{r.taskId}: TaskId; \texttt{r.requester}: Address;
  \texttt{r.payment}: $\mathbb{N}$ -- the escrowed payment offered for the
  request
\item \texttt{r.status}: RequestStatus
\item \texttt{r.minProviders}, \texttt{r.maxProviders}: $\mathbb{N}$ --
  copied from the task's parameters when the request is created
\item \texttt{r.createdAt}, \texttt{r.acceptanceDeadline},
  \texttt{r.acceptanceEnd}, \texttt{r.commitDeadline}, \texttt{r.commitEnd},
  \texttt{r.revealDeadline}, \texttt{r.revealEnd},
  \texttt{r.pendingDeadline}: Timestamp (or $\bot$ until set)
\item \texttt{r.accepted}: sequence of Address, in order of acceptance
\item \texttt{r.commitment}: Address $\rightarrow$ Hash $\cup$ \{$\bot$\}
\item \texttt{r.prediction}: Address $\rightarrow$ $\mathbb{R}$ $\cup$
  \{$\bot$\} -- set only for valid reveals
\item \texttt{r.outcome}: Address $\rightarrow$ ParticipationOutcome -- the
  latest participation outcome of each accepted provider
\item \texttt{r.snapshot}: Address $\rightarrow$ $\mathbb{R}$ -- each
  provider's reputation when the reveal phase ends
\item \texttt{r.result}: $\mathbb{R}$ $\cup$ \{$\bot$\};
  \texttt{r.lowest}, \texttt{r.highest}: $\mathbb{R}$ $\cup$ \{$\bot$\};
  \texttt{r.included}: $\mathbb{N}$
\item \texttt{r.unresolved}: Address $\rightarrow$ $\mathbb{N}$ -- rewards
  held in escrow because the transfer failed
\end{itemize}

\subsubsection{Environment Variables}

\begin{itemize}
\item \texttt{ChainClock}: Timestamp -- the blockchain's current time for the
  executing transaction
\end{itemize}

\subsubsection{Assumptions}

\begin{itemize}
\item The \texttt{Address} arguments of \texttt{createRequest},
  \texttt{acceptRequest}, \texttt{submitCommitment}, \texttt{revealPrediction},
  and \texttt{retryTransfer} are the signers of the transaction, so a caller
  can act only for its own address.
\item The on-chain program does not run on its own. \texttt{advancePhase} is
  invoked by any party, and the request advances only when it is called.
\item The hash function behind \texttt{computeCommitment} is collision- and
  preimage-resistant. Salts are at least 256 bits and are generated by the
  Participation Client (\mref{mParticipation}).
\item A transaction that raises an exception leaves all state unchanged.
\item The Provider Registry's recorded stake for a provider and task equals the
  balance of that provider's stake escrow in the Token Manager. Modules that
  move stake tokens (Penalty Enforcer, Reputation Handler) leave the Registry's
  record unchanged, and this module updates the record afterwards.
\item Refunds to the requester and payments to providers with a valid reveal
  always succeed, because those wallets were able to pay in or to reveal.
\item The other modules behave as specified in their own MIS.
\end{itemize}

\subsubsection{Access Routine Semantics}

\noindent \texttt{computeCommitment}(id, provider, p, salt):
\begin{itemize}
\item output: the hash of the concatenation of \texttt{id}, \texttt{provider},
  \texttt{p}, and \texttt{salt}, in that order.
\end{itemize}

\noindent \texttt{createRequest}(taskId, requester, payment):
\begin{itemize}
\item transition: Let \texttt{params} := \texttt{getTaskParameters(taskId)}.
  Allocate the identifier \texttt{id}. Call \texttt{escrowPayment(id,
  requester, payment)}. Create a request \texttt{r} with \texttt{r.status} :=
  Acceptance, \texttt{r.createdAt} := \texttt{ChainClock},
  \texttt{r.acceptanceDeadline} := \texttt{ChainClock} +
  \texttt{params.acceptancePeriod}, and \texttt{minProviders} and
  \texttt{maxProviders} copied from \texttt{params}. All other fields start
  empty.
\item output: the identifier \texttt{id} of the new request.
\item exception: \texttt{UnknownTaskError} if the task does not exist.
  \texttt{InvalidPaymentError} if \texttt{payment} $\le$
  \texttt{MINIMUM\_REQUEST\_PAYMENT}. \texttt{TransferFailedError} if the
  requester's wallet cannot cover the payment.
\end{itemize}

\noindent \texttt{acceptRequest}(id, provider):
\begin{itemize}
\item transition: Let \texttt{reservation} := (\texttt{r.minProviders} $-$ 1)
  $\times$ \texttt{r.payment}. Call \texttt{reserveStake(r.taskId, provider,
  reservation)}. Append \texttt{provider} to \texttt{r.accepted} and set
  \texttt{r.outcome[provider]} := NoCommit. If $|\texttt{r.accepted}| =
  \texttt{r.maxProviders}$, close the acceptance period early:
  \texttt{r.acceptanceEnd} := \texttt{ChainClock}, \texttt{r.status} := Commit,
  and \texttt{r.commitDeadline} := \texttt{r.acceptanceEnd} +
  \texttt{COMMIT\_PERIOD}.
\item exception: \texttt{WrongPhaseError} if \texttt{r.status} $\neq$
  Acceptance. \texttt{DeadlinePassedError} if \texttt{ChainClock} $\ge$
  \texttt{r.acceptanceDeadline}. \texttt{NotRegisteredError} if
  \texttt{isRegistered(r.taskId, provider)} is false.
  \texttt{AlreadyAcceptedError} if \texttt{provider} $\in$
  \texttt{r.accepted}. \texttt{InsufficientStakeError} if
  \texttt{getStake(r.taskId, provider)} $-$ \texttt{getReservedStake(r.taskId,
  provider)} $-$ \texttt{reservation} $<$ \texttt{MINIMUM\_PROVIDER\_STAKE}.
\end{itemize}

\noindent \texttt{submitCommitment}(id, provider, c):
\begin{itemize}
\item transition: \texttt{r.commitment[provider]} := \texttt{c} and
  \texttt{r.outcome[provider]} := NoReveal. If every provider in
  \texttt{r.accepted} now has a commitment, then \texttt{r.commitEnd} :=
  \texttt{ChainClock}, \texttt{r.status} := Reveal, and
  \texttt{r.revealDeadline} := \texttt{r.commitEnd} + \texttt{REVEAL\_PERIOD}.
\item exception: \texttt{WrongPhaseError} if \texttt{r.status} $\neq$ Commit.
  \texttt{DeadlinePassedError} if \texttt{ChainClock} $\ge$
  \texttt{r.commitDeadline}. \texttt{NotAcceptedError} if \texttt{provider}
  $\notin$ \texttt{r.accepted}. \texttt{AlreadyCommittedError} if
  \texttt{r.commitment[provider]} $\neq$ $\bot$.
\end{itemize}

\noindent \texttt{revealPrediction}(id, provider, p, salt):
\begin{itemize}
\item transition: Let \texttt{matches} := (\texttt{computeCommitment(id,
  provider, p, salt)} = \texttt{r.commitment[provider]}) and
  \texttt{conforms} := \texttt{isValidOutput(r.taskId, p)}.
  \begin{itemize}
  \item If \texttt{matches} $\wedge$ \texttt{conforms}:
    \texttt{r.prediction[provider]} := \texttt{p} and
    \texttt{r.outcome[provider]} := ValidReveal.
  \item If $\neg$\texttt{matches}: \texttt{r.outcome[provider]} :=
    RevealMismatch.
  \item Otherwise: \texttt{r.outcome[provider]} := RevealNonConforming.
  \end{itemize}
  In every case the provider's single reveal is consumed. If every provider
  with a commitment has now revealed, then \texttt{r.revealEnd} :=
  \texttt{ChainClock}, \texttt{r.snapshot} := the current reputations of the
  committed providers, and \texttt{r.status} := Finalization.
\item output: Accepted, Mismatch, or NonConforming, matching the three cases.
  A rejected reveal is reported through the output, not an exception, because
  an exception would roll back the record of the offence. The offence is
  penalized when the request ends.
\item exception: \texttt{WrongPhaseError} if \texttt{r.status} $\neq$ Reveal.
  \texttt{DeadlinePassedError} if \texttt{ChainClock} $\ge$
  \texttt{r.revealDeadline}. \texttt{NoCommitmentError} if
  \texttt{r.commitment[provider]} = $\bot$. \texttt{AlreadyRevealedError} if
  \texttt{r.outcome[provider]} $\neq$ NoReveal.
\end{itemize}

\noindent \texttt{advancePhase}(id):
\begin{itemize}
\item transition: Repeat the first applicable rule below until none applies.
  \begin{itemize}
  \item \emph{Acceptance} and \texttt{ChainClock} $\ge$
    \texttt{r.acceptanceDeadline}: \texttt{r.acceptanceEnd} :=
    \texttt{r.acceptanceDeadline}. If $|\texttt{r.accepted}| <
    \texttt{r.minProviders}$, run \texttt{cancel(r)}. Otherwise
    \texttt{r.status} := Commit and \texttt{r.commitDeadline} :=
    \texttt{r.acceptanceEnd} + \texttt{COMMIT\_PERIOD}.
  \item \emph{Commit} and \texttt{ChainClock} $\ge$ \texttt{r.commitDeadline}:
    \texttt{r.commitEnd} := \texttt{r.commitDeadline}. If no provider has a
    commitment, then \texttt{r.revealEnd} := \texttt{r.commitEnd},
    \texttt{r.snapshot} := empty, and \texttt{r.status} := Finalization.
    Otherwise \texttt{r.status} := Reveal and \texttt{r.revealDeadline} :=
    \texttt{r.commitEnd} + \texttt{REVEAL\_PERIOD}.
  \item \emph{Reveal} and \texttt{ChainClock} $\ge$ \texttt{r.revealDeadline}:
    \texttt{r.revealEnd} := \texttt{r.revealDeadline}, \texttt{r.snapshot} :=
    the current reputations of the committed providers, and
    \texttt{r.status} := Finalization.
  \item \emph{Finalization}: let $I$ be the set of providers whose outcome is
    ValidReveal and whose snapshot reputation is non-zero. If $|I| <
    \texttt{r.minProviders}$, run \texttt{cancel(r)}. Otherwise run
    \texttt{finalize(r, I)}.
  \item \emph{PendingTransfer} and \texttt{ChainClock} $\ge$
    \texttt{r.pendingDeadline}: for each unresolved reward, call
    \texttt{forfeitReward(id, amount)} and discard the entry. Then
    \texttt{r.status} := Completed.
  \end{itemize}
\item output: \texttt{r.status} after the rules have been applied.
\item exception: \texttt{UnknownRequestError} if \texttt{id} does not exist.
  \texttt{ArithmeticOverflowError} if an overflow occurs in a called module;
  the request is then left unchanged.
\end{itemize}

\noindent \texttt{retryTransfer}(id, provider):
\begin{itemize}
\item transition: Call \texttt{retryReward(id, provider,
  r.unresolved[provider])}, then remove \texttt{provider} from
  \texttt{r.unresolved}. If \texttt{r.unresolved} is now empty,
  \texttt{r.status} := Completed.
\item exception: \texttt{WrongPhaseError} if \texttt{r.status} $\neq$
  PendingTransfer. \texttt{DeadlinePassedError} if \texttt{ChainClock} $\ge$
  \texttt{r.pendingDeadline}. \texttt{NoUnresolvedTransferError} if
  \texttt{provider} $\notin$ \texttt{r.unresolved}.
  \texttt{TransferFailedError} if the transfer fails again.
\end{itemize}

\noindent \texttt{getStatus}(id), \texttt{getResult}(id),
\texttt{getDeadlines}(id), \texttt{getUnresolvedTransfers}(id):
\begin{itemize}
\item output: \texttt{r.status}; (\texttt{r.result}, \texttt{r.lowest},
  \texttt{r.highest}, \texttt{r.included}); the tuple
  (\texttt{r.acceptanceDeadline}, \texttt{r.commitDeadline},
  \texttt{r.revealDeadline}, \texttt{r.pendingDeadline}); and
  \texttt{r.unresolved} as a set of pairs.
\item exception: \texttt{UnknownRequestError} for every routine.
  \texttt{ResultUnavailableError} for \texttt{getResult} if \texttt{r.result}
  = $\bot$.
\end{itemize}

\subsubsection{Local Functions}

\begin{itemize}
\item \texttt{releaseReservations(r)}: calls \texttt{releaseStake} for each
  provider in \texttt{r.accepted}, releasing the reservation made in
  \texttt{acceptRequest}, less any amount already collected as compensation.
\item \texttt{applyPenalties(r)}: for each provider \texttt{a} in
  \texttt{r.accepted} that \texttt{isRegistered(r.taskId, a)}, let
  \texttt{d} := \texttt{applyPenalty(r.taskId, a, r.outcome[a],
  getStake(r.taskId, a), getSlashHistory(r.taskId, a))}. This burns the
  slashed tokens, and also the remaining stake if \texttt{d.burnRemaining}.
  Then update the Registry's record: if \texttt{d.unregister}, call
  \texttt{removeProvider(r.taskId, a)}; otherwise, if \texttt{d.slashAmount}
  $> 0$, call \texttt{deductStake(r.taskId, a, d.slashAmount)}. Finally call
  \texttt{recordParticipation(r.taskId, a, id, d.slashAmount} $> 0$\texttt{)}.
\item \texttt{cancel(r)}: let $V$ be the providers whose outcome is
  ValidReveal and $F$ the remaining providers in \texttt{r.accepted}, in
  acceptance order. If $V \neq \emptyset$, let $C$ := \texttt{r.payment}
  $\times |V|$, call \texttt{collectCompensation(r.taskId, id, F, C)}, call
  \texttt{deductStake(r.taskId, f, share)} for each provider $f$ in $F$ with
  the share returned for it, and call \texttt{payCompensation(id, V,
  r.payment)}. After that, call \texttt{refundPayment(id, r.requester,
  r.payment)}, run \texttt{releaseReservations(r)} and
  \texttt{applyPenalties(r)}, and set \texttt{r.status} := Cancelled.
\item \texttt{finalize(r, I)}: let $p$ be the predictions of $I$ and $\rho$
  their snapshot reputations, in acceptance order.
  \begin{enumerate}
  \item Let \texttt{F} := \texttt{weightedMedian(p, $\rho$)} and
    (\texttt{lowest}, \texttt{highest}) := \texttt{resultSpread(p)}. If
    \texttt{AggregationError} is raised, call \texttt{refundPayment(id,
    r.requester, r.payment)}, run \texttt{releaseReservations(r)} and
    \texttt{applyPenalties(r)}, set \texttt{r.status} := Failed, and stop.
  \item Let $\rho'$ := \texttt{updateReputations(p, $\rho$, F)} and call
    \texttt{setReputation} for each provider in $I$ with its new value. For
    each provider $a$ in $I$ whose new reputation is 0, call
    \texttt{forfeitStake(r.taskId, a, getStake(r.taskId, a))} and then
    \texttt{removeProvider(r.taskId, a)}.
  \item Set \texttt{r.unresolved} := \texttt{distributeRewards(id,
    r.requester, r.payment, I, $\rho$, $\rho'$)}.
  \item Set \texttt{r.result} := \texttt{F}, \texttt{r.lowest},
    \texttt{r.highest}, and \texttt{r.included} := $|I|$. Run
    \texttt{releaseReservations(r)} and \texttt{applyPenalties(r)}.
  \item Set \texttt{r.pendingDeadline} := \texttt{r.revealEnd} +
    \texttt{PENDING\_TRANSFER\_PERIOD}. Set \texttt{r.status} := Completed if
    \texttt{r.unresolved} is empty, and PendingTransfer otherwise.
  \end{enumerate}
\end{itemize}


% =====================================================================
\section{MIS of Prediction Aggregator (\mref{mAggregator})} \label{misAggregator}
% =====================================================================

\subsection{Module}

Prediction Aggregator (\mref{mAggregator}). A library implemented in the
on-chain program.

\subsection{Uses}

This module uses no other module.

\subsection{Syntax}

\subsubsection{Exported Constants}

None.

\subsubsection{Exported Access Programs}

\begin{center}
\begin{tabular}{p{3cm} p{4.5cm} p{3cm} p{5cm}}
\hline
\textbf{Name} & \textbf{In} & \textbf{Out} & \textbf{Exceptions} \\
\hline
weightedMedian & seq of $\mathbb{R}$ (predictions), seq of $\mathbb{R}$
  (reputations) & $\mathbb{R}$ & EmptyInputError, LengthMismatchError,
  InvalidReputationError, ArithmeticOverflowError, AggregationError \\
\hline
resultSpread & seq of $\mathbb{R}$ (predictions) & $\mathbb{R} \times
  \mathbb{R}$ & EmptyInputError \\
\hline
\end{tabular}
\end{center}

\subsection{Semantics}

\subsubsection{State Variables}

None.

\subsubsection{Environment Variables}

None.

\subsubsection{Assumptions}

\begin{itemize}
\item The caller passes only valid revealed predictions, each paired by
  position with the reputation of its provider at the end of the reveal phase,
  and only from providers with non-zero reputation.
\item The caller checks that the number of predictions meets the task's
  minimum before calling.
\item Implementations compute with exact integer or fixed-point arithmetic, so
  the result equals the mathematical definition below whenever no exception is
  raised.
\end{itemize}

\subsubsection{Access Routine Semantics}

\noindent \texttt{weightedMedian}($P$, $R$):
\begin{itemize}
\item output: Let $n = |P|$ and let $\pi$ be a permutation that sorts $P$ in
  nondecreasing order (equal predictions keep their input order). Let $S =
  \sum_{i} R_i$ and $S_k = \sum_{j=1}^{k} R_{\pi(j)}$, with $S_0 = 0$. Let
  $k^*$ be the smallest $k$ such that $2 S_k > S$. The output is
  \[
  \begin{cases}
  \dfrac{P_{\pi(k^*-1)} + P_{\pi(k^*)}}{2} & \text{if } 2 S_{k^*-1} = S,\\[1ex]
  P_{\pi(k^*)} & \text{otherwise.}
  \end{cases}
  \]
  This is the weighted median of Appendix~A of the SRS, where each prediction's
  weight is $R_i / S$. The comparisons use $2 S_k$ and $S$ directly, so no
  division is needed.
\item exception: \texttt{EmptyInputError} if $n = 0$.
  \texttt{LengthMismatchError} if $|P| \neq |R|$.
  \texttt{InvalidReputationError} if some $R_i \le 0$.
  \texttt{ArithmeticOverflowError} if any intermediate value is outside the
  representable range. \texttt{AggregationError} if no $k^*$ exists, which
  cannot happen when the assumptions hold.
\end{itemize}

\noindent \texttt{resultSpread}($P$):
\begin{itemize}
\item output: $(\min_i P_i, \max_i P_i)$.
\item exception: \texttt{EmptyInputError} if $P$ is empty.
\end{itemize}

\subsubsection{Local Functions}

\begin{itemize}
\item \texttt{sortedOrder}($P$): the permutation $\pi$ used above.
\item \texttt{cumulativeReputation}($R$, $\pi$, $k$): the sum $S_k$.
\end{itemize}


% =====================================================================
\section{MIS of Rewards Calculator (\mref{mRewards})} \label{misRewards}
% =====================================================================

\subsection{Module}

Rewards Calculator (\mref{mRewards}). A library implemented in the on-chain
program. It decides who earns a reward and how much, and carries out every
payment from and refund of a request's reward escrow. The escrow balances
themselves are held by the Token Manager, so this module has no state of its
own.

\subsection{Uses}

\begin{itemize}
\item \textbf{Token Manager (\mref{mToken}):}
  \texttt{transfer(source, destination, amount)} with the endpoints
  \texttt{Wallet(a)}, \texttt{RequestEscrow(id)}, and \texttt{Burn}, and the
  exception \texttt{TransferFailedError}.
\end{itemize}

\subsection{Syntax}

\subsubsection{Exported Constants}

None.

\subsubsection{Exported Access Programs}

\begin{longtable}{p{3.2cm} p{5.3cm} p{3cm} p{4.1cm}}
\hline
\textbf{Name} & \textbf{In} & \textbf{Out} & \textbf{Exceptions} \\
\hline
\endfirsthead
\hline
\textbf{Name} & \textbf{In} & \textbf{Out} & \textbf{Exceptions} \\
\hline
\endhead
eligibleProviders & seq of $\mathbb{R}$ (reputations before), seq of
  $\mathbb{R}$ (reputations after) & seq of $\mathbb{B}$ &
  EmptyInputError, LengthMismatchError \\
\hline
calculateRewards & $\mathbb{N}$ (total reward), seq of $\mathbb{R}$
  (reputations before), seq of $\mathbb{R}$ (reputations after) & seq of
  $\mathbb{N}$ (shares) $\times$ $\mathbb{N}$ (remainder) & EmptyInputError,
  LengthMismatchError, NoEligibleProvidersError, ArithmeticOverflowError,
  RewardValidationError \\
\hline
isValidDistribution & $\mathbb{N}$ (total reward), seq of $\mathbb{N}$
  (shares) & $\mathbb{B}$ & - \\
\hline
escrowPayment & RequestId, Address (requester), $\mathbb{N}$ & - &
  TransferFailedError \\
\hline
distributeRewards & RequestId, Address (requester), $\mathbb{N}$ (total
  reward), seq of Address, seq of $\mathbb{R}$ (reputations before), seq of
  $\mathbb{R}$ (reputations after) & seq of (Address $\times$ $\mathbb{N}$)
  (unresolved rewards) & the exceptions of calculateRewards \\
\hline
retryReward & RequestId, Address, $\mathbb{N}$ & - & TransferFailedError \\
\hline
forfeitReward & RequestId, $\mathbb{N}$ & - & TransferFailedError \\
\hline
refundPayment & RequestId, Address (requester), $\mathbb{N}$ & - &
  TransferFailedError \\
\hline
payCompensation & RequestId, seq of Address, $\mathbb{N}$ (amount each) & - &
  TransferFailedError \\
\hline
\end{longtable}

\subsection{Semantics}

\subsubsection{State Variables}

None.

\subsubsection{Environment Variables}

None.

\subsubsection{Assumptions}

\begin{itemize}
\item The sequences are aligned by position and contain only the providers
  whose valid predictions were included in the aggregate. Each ``before'' value
  is the reputation snapshot from the end of the reveal phase and each
  ``after'' value is the result of the reputation update for the same request.
\item Reputations are exact values, so products and quotients are computed
  without rounding until the final rounding down of each share.
\item Each request's escrow in the Token Manager holds at least the amounts
  this module is asked to pay out.
\item Refunds to the requester succeed, because the requester's wallet was able
  to pay into escrow.
\end{itemize}

\subsubsection{Access Routine Semantics}

\noindent \texttt{eligibleProviders}($R$, $R'$):
\begin{itemize}
\item output: the sequence $E$ with $E_i = \top$ if $R'_i \ge R_i$ (the
  provider's reputation did not decrease) and $\bot$ otherwise.
\item exception: \texttt{EmptyInputError} if $R$ is empty.
  \texttt{LengthMismatchError} if $|R| \neq |R'|$.
\end{itemize}

\noindent \texttt{calculateRewards}($T$, $R$, $R'$):
\begin{itemize}
\item output: Let $E$ := \texttt{eligibleProviders}($R$, $R'$) and $W =
  \sum_{i : E_i} R'_i$. The shares are $s_i = \left\lfloor T \cdot R'_i / W
  \right\rfloor$ if $E_i$, and $s_i = 0$ otherwise. The remainder is $T -
  \sum_i s_i$. Shares of eligible providers are proportional to their
  reputation after the update.
\item exception: \texttt{NoEligibleProvidersError} if no provider is eligible.
  \texttt{ArithmeticOverflowError} if $T \cdot R'_i$ or any sum is outside the
  representable range. \texttt{RewardValidationError} if
  \texttt{isValidDistribution}($T$, $s$) is false. A provider whose
  calculated share is zero simply receives zero. The other exceptions are
  those of \texttt{eligibleProviders}.
\end{itemize}

\noindent \texttt{isValidDistribution}($T$, $s$):
\begin{itemize}
\item output: $\top$ if every $s_i \ge 0$ and $\sum_i s_i \le T$, and $\bot$
  otherwise.
\end{itemize}

\noindent \texttt{escrowPayment}(id, requester, amount):
\begin{itemize}
\item transition: \texttt{transfer(Wallet(requester), RequestEscrow(id),
  amount)}.
\item exception: \texttt{TransferFailedError} if the requester's wallet cannot
  cover \texttt{amount}.
\end{itemize}

\noindent \texttt{distributeRewards}(id, requester, total, providers, $R$, $R'$):
\begin{itemize}
\item transition: Let ($s$, \texttt{remainder}) :=
  \texttt{calculateRewards}(total, $R$, $R'$). For each index $i$ with $s_i >
  0$, attempt \texttt{transfer(RequestEscrow(id), Wallet(providers$_i$),
  $s_i$)}. A \texttt{TransferFailedError} on one provider does not stop the
  others; the provider's reward stays in escrow and is added to the output.
  Finally \texttt{transfer(RequestEscrow(id), Wallet(requester),
  remainder)}.
\item output: the pairs (\texttt{providers}$_i$, $s_i$) whose transfer failed.
\item exception: those of \texttt{calculateRewards}.
\end{itemize}

\noindent \texttt{retryReward}(id, provider, amount):
\begin{itemize}
\item transition: \texttt{transfer(RequestEscrow(id), Wallet(provider),
  amount)}.
\item exception: \texttt{TransferFailedError} if the provider's wallet still
  cannot receive tokens.
\end{itemize}

\noindent \texttt{forfeitReward}(id, amount):
\begin{itemize}
\item transition: \texttt{transfer(RequestEscrow(id), Burn, amount)}.
\item exception: \texttt{TransferFailedError} if the escrow holds less than
  \texttt{amount}.
\end{itemize}

\noindent \texttt{refundPayment}(id, requester, amount):
\begin{itemize}
\item transition: \texttt{transfer(RequestEscrow(id), Wallet(requester),
  amount)}.
\item exception: \texttt{TransferFailedError} if the escrow holds less than
  \texttt{amount}.
\end{itemize}

\noindent \texttt{payCompensation}(id, providers, amountEach):
\begin{itemize}
\item transition: for each provider, \texttt{transfer(RequestEscrow(id),
  Wallet(provider), amountEach)}. The escrow has been topped up beforehand by
  Penalty Enforcer's \texttt{collectCompensation}.
\item exception: \texttt{TransferFailedError} if a transfer fails; no payment
  is then made.
\end{itemize}

\subsubsection{Local Functions}

\begin{itemize}
\item \texttt{eligibleReputationTotal}($R$, $R'$): the sum $W$ used above.
\end{itemize}


% =====================================================================
\section{MIS of Penalty Enforcer (\mref{mPenalty})} \label{misPenalty}
% =====================================================================

\subsection{Module}

Penalty Enforcer (\mref{mPenalty}). A library implemented in the on-chain
program. It decides whether a penalty applies and how large it is, and it burns
the slashed stake or collects it as compensation. Stake balances are held by
the Token Manager, so this module has no state of its own.

\subsection{Uses}

\begin{itemize}
\item \textbf{Token Manager (\mref{mToken}):}
  \texttt{transfer(source, destination, amount)} with the endpoints
  \texttt{StakeEscrow(task, a)}, \texttt{RequestEscrow(id)}, and
  \texttt{Burn}, and the exception \texttt{TransferFailedError}.
\end{itemize}

\subsection{Syntax}

\subsubsection{Exported Constants}

\begin{itemize}
\item \texttt{SLASH\_FRACTION}: $\mathbb{R}$ -- the fraction of a provider's
  stake for a task that is slashed for an offence, with $0 < $
  \texttt{SLASH\_FRACTION} $\le 1$
\item \texttt{REPEAT\_OFFENSE\_COUNT}: $\mathbb{N}$ -- the number of slashed
  requests, within the window, that makes a provider a repeat offender
\item \texttt{REPEAT\_OFFENSE\_WINDOW}: $\mathbb{N}$ -- the number of most recent
  accepted requests over which offences are counted, with $1 \le$
  \texttt{REPEAT\_OFFENSE\_COUNT} $\le$ \texttt{REPEAT\_OFFENSE\_WINDOW}
\end{itemize}

\subsubsection{Exported Types}
\begin{itemize}
\item \texttt{ParticipationOutcome} = \{NoCommit, NoReveal, RevealMismatch,
  RevealNonConforming, ValidReveal\}
\item \texttt{PenaltyDecision} = (\texttt{slashAmount}: $\mathbb{N}$,
  \texttt{unregister}: $\mathbb{B}$, \texttt{burnRemaining}: $\mathbb{B}$)
\end{itemize}

\subsubsection{Exported Access Programs}

\begin{longtable}{p{3.2cm} p{5.3cm} p{3cm} p{4.1cm}}
\hline
\textbf{Name} & \textbf{In} & \textbf{Out} & \textbf{Exceptions} \\
\hline
\endfirsthead
\hline
\textbf{Name} & \textbf{In} & \textbf{Out} & \textbf{Exceptions} \\
\hline
\endhead
isOffense & ParticipationOutcome & $\mathbb{B}$ & - \\
\hline
slashAmount & $\mathbb{N}$ (stake) & $\mathbb{N}$ & ArithmeticOverflowError \\
\hline
penaltyFor & ParticipationOutcome, $\mathbb{N}$ (stake), seq of
  $\mathbb{B}$ (slash history) & PenaltyDecision & ArithmeticOverflowError \\
\hline
compensationShares & $\mathbb{N}$ (total), $\mathbb{N}$ (number of
  providers) & seq of $\mathbb{N}$ & EmptyInputError \\
\hline
applyPenalty & TaskId, Address, ParticipationOutcome, $\mathbb{N}$ (stake),
  seq of $\mathbb{B}$ (slash history) & PenaltyDecision &
  ArithmeticOverflowError, TransferFailedError \\
\hline
collectCompensation & TaskId, RequestId, seq of Address, $\mathbb{N}$ (total)
  & seq of $\mathbb{N}$ (share taken from each) & EmptyInputError,
  ArithmeticOverflowError, TransferFailedError \\
\hline
\end{longtable}

\subsection{Semantics}

\subsubsection{State Variables}

None.

\subsubsection{Environment Variables}

None.

\subsubsection{Assumptions}

\begin{itemize}
\item The slash history lists, for the provider's most recent accepted requests
  for the task before the current one, whether the provider was slashed
  ($\top$) or not ($\bot$), with the most recent request last. Entries beyond
  the most recent \texttt{REPEAT\_OFFENSE\_WINDOW} $-$ 1 are ignored.
\item The stake passed in is the provider's current stake for the task, equal
  to the balance of \texttt{StakeEscrow(task, a)}.
\item The caller updates the Provider Registry's stake record after a call
  that moves stake tokens.
\item The constants satisfy the ranges given above.
\end{itemize}

\subsubsection{Access Routine Semantics}

\noindent \texttt{isOffense}(o):
\begin{itemize}
\item output: $\top$ if $o \in$ \{NoCommit, NoReveal, RevealMismatch,
  RevealNonConforming\}, and $\bot$ if $o$ = ValidReveal.
\end{itemize}

\noindent \texttt{slashAmount}(stake):
\begin{itemize}
\item output: $\lfloor$ \texttt{SLASH\_FRACTION} $\times$ stake $\rfloor$.
\item exception: \texttt{ArithmeticOverflowError} if the product is outside
  the representable range.
\end{itemize}

\noindent \texttt{penaltyFor}(o, stake, history):
\begin{itemize}
\item output: Let $\mathit{off}$ := \texttt{isOffense}(o), let $h$ be the last
  \texttt{REPEAT\_OFFENSE\_WINDOW} $-$ 1 entries of \texttt{history}, and let
  $c = [\mathit{off}] + |\{x \in h : x = \top\}|$, where $[\mathit{off}]$ is 1
  if $\mathit{off}$ and 0 otherwise. Let $\mathit{repeat} = \mathit{off}
  \wedge c \ge$ \texttt{REPEAT\_OFFENSE\_COUNT}. The output is
  (\texttt{slashAmount}: \texttt{slashAmount}(stake) if $\mathit{off}$, else 0;
  \texttt{unregister}: $\mathit{repeat}$; \texttt{burnRemaining}:
  $\mathit{repeat}$).
\item exception: \texttt{ArithmeticOverflowError}, as for
  \texttt{slashAmount}.
\end{itemize}

\noindent \texttt{compensationShares}(total, n):
\begin{itemize}
\item output: the sequence of length $n$ whose $i$-th entry (counting from 1)
  is $\lfloor \mathit{total} / n \rfloor$ plus 1 if $i \le \mathit{total} \bmod
  n$. The entries sum to \texttt{total}, and the earliest providers carry the
  remainder.
\item exception: \texttt{EmptyInputError} if $n = 0$.
\end{itemize}

\noindent \texttt{applyPenalty}(task, a, o, stake, history):
\begin{itemize}
\item transition: Let $d$ := \texttt{penaltyFor}(o, stake, history). If
  $d$.\texttt{slashAmount} $> 0$, perform \texttt{transfer(StakeEscrow(task,
  a), Burn, d.slashAmount)}. If $d$.\texttt{burnRemaining} and \texttt{stake}
  $-$ $d$.\texttt{slashAmount} $> 0$, perform \texttt{transfer(StakeEscrow(task,
  a), Burn, stake} $-$ $d$\texttt{.slashAmount)}.
\item output: $d$.
\item exception: those of \texttt{penaltyFor}, and \texttt{TransferFailedError}
  if a transfer fails.
\end{itemize}

\noindent \texttt{collectCompensation}(task, id, providers, total):
\begin{itemize}
\item transition: Let \texttt{shares} := \texttt{compensationShares}(total,
  $|$providers$|$). For each index $i$, perform
  \texttt{transfer(StakeEscrow(task, providers$_i$), RequestEscrow(id),
  shares$_i$)}.
\item output: \texttt{shares}.
\item exception: \texttt{EmptyInputError} if \texttt{providers} is empty.
  \texttt{ArithmeticOverflowError} as above. \texttt{TransferFailedError} if
  a provider's stake escrow holds less than its share; no transfer is then
  made.
\end{itemize}

\subsubsection{Local Functions}

\begin{itemize}
\item \texttt{recentWindow}(history): the last
  \texttt{REPEAT\_OFFENSE\_WINDOW} $-$ 1 entries of \texttt{history}.
\end{itemize}


% =====================================================================
\section{MIS of Reputation Handler (\mref{mReputation})} \label{misReputation}
% =====================================================================

\subsection{Module}

Reputation Handler (\mref{mReputation}). A library implemented in the on-chain
program. It calculates reputation updates, and burns the stake of a provider
whose reputation reaches 0.

\subsection{Uses}

\begin{itemize}
\item \textbf{Token Manager (\mref{mToken}):}
  \texttt{transfer(source, destination, amount)} with the endpoints
  \texttt{StakeEscrow(task, a)} and \texttt{Burn}, and the exception
  \texttt{TransferFailedError}.
\end{itemize}

\subsection{Syntax}

\subsubsection{Exported Constants}

\begin{itemize}
\item \texttt{REPUTATION\_MIN}: $\mathbb{R}$ = 0
\item \texttt{REPUTATION\_MAX}: $\mathbb{R}$ = 100
\item \texttt{REPUTATION\_ADJUSTMENT\_FACTOR}: $\mathbb{R}$ = 1.0 -- the factor
  $\alpha$ that scales each reputation change
\end{itemize}

\subsubsection{Exported Access Programs}

\begin{center}
\begin{tabular}{p{3cm} p{5.2cm} p{3cm} p{4.3cm}}
\hline
\textbf{Name} & \textbf{In} & \textbf{Out} & \textbf{Exceptions} \\
\hline
updateReputations & seq of $\mathbb{R}$ (predictions), seq of $\mathbb{R}$
  (reputations), $\mathbb{R}$ (aggregate) & seq of $\mathbb{R}$ &
  EmptyInputError, LengthMismatchError, InvalidReputationError,
  ArithmeticOverflowError \\
\hline
clampReputation & $\mathbb{R}$ & $\mathbb{R}$ & - \\
\hline
forfeitStake & TaskId, Address, $\mathbb{N}$ (stake) & - &
  TransferFailedError \\
\hline
\end{tabular}
\end{center}

\subsection{Semantics}

\subsubsection{State Variables}

None.

\subsubsection{Environment Variables}

None.

\subsubsection{Assumptions}

\begin{itemize}
\item The reputations passed to \texttt{updateReputations} are the snapshot
  from the end of the reveal phase for providers with a valid prediction
  included in the aggregate, aligned by position with the predictions.
\item The aggregate is the weighted median that the Prediction Aggregator
  (\mref{mAggregator}) produced from the same predictions and reputations.
\item \texttt{forfeitStake} is called only for a provider whose new reputation
  is 0. The caller removes the provider from the Provider Registry afterwards.
\end{itemize}

\subsubsection{Access Routine Semantics}

\noindent \texttt{clampReputation}(x):
\begin{itemize}
\item output: $\min($\texttt{REPUTATION\_MAX}$, \max($\texttt{REPUTATION\_MIN}$,
  x))$.
\end{itemize}

\noindent \texttt{updateReputations}($P$, $R$, $F$):
\begin{itemize}
\item output: Let $\sigma^2 = \sum_i R_i (P_i - F)^2 / \sum_i R_i$. If
  $\sigma^2 = 0$, every reputation change is 0 and the output is
  \texttt{clampReputation}($R_i$) for each $i$. Otherwise let $\sigma =
  \sqrt{\sigma^2}$ and, for each $i$, let $\delta_i =
  \alpha\left(1 - |F - P_i| / \sigma\right)$ with $\alpha$ =
  \texttt{REPUTATION\_ADJUSTMENT\_FACTOR}. The output is
  \texttt{clampReputation}($R_i + \delta_i$) for each $i$. This is the
  confidence-based update of Appendix~A of the SRS. A provider closer to $F$
  than one standard deviation gains reputation and one farther loses it, and
  all values stay within [\texttt{REPUTATION\_MIN}, \texttt{REPUTATION\_MAX}].
\item exception: \texttt{EmptyInputError} if $P$ is empty.
  \texttt{LengthMismatchError} if $|P| \neq |R|$.
  \texttt{InvalidReputationError} if some $R_i$ is not in
  (\texttt{REPUTATION\_MIN}, \texttt{REPUTATION\_MAX}].
  \texttt{ArithmeticOverflowError} if any intermediate value (including the
  squared differences and the square root) is outside the representable range.
\end{itemize}

\noindent \texttt{forfeitStake}(task, a, stake):
\begin{itemize}
\item transition: if \texttt{stake} $> 0$, perform
  \texttt{transfer(StakeEscrow(task, a), Burn, stake)}.
\item exception: \texttt{TransferFailedError} if the stake escrow holds less
  than \texttt{stake}.
\end{itemize}

\subsubsection{Local Functions}

\begin{itemize}
\item \texttt{weightedVariance}($P$, $R$, $F$): the value $\sigma^2$ above.
\end{itemize}



\newpage

\bibliographystyle {plainnat}
\bibliography {../../../refs/References}

\newpage

\newpage{}

\section{Appendix --- Generative AI Use Disclosure}

\wss{ChatGPT (version 5) was used to brainstorm ideas for the template for this disclose section.}

\subsection{AI Tools Used}

\wss{Give the name of the AI tool(s) used and the version.}

\subsection{How and Where Used}

\wss{Explain what the tool(s) were used for.  Examples include: brainstorming,
explaining a technical concept, reviewing the document, generating or modifying
text, generating or modifying diagrams, identifying errors or omissions or
inconsistencies.}

\wss{For each of the uses, identify which parts of the document were affected.}

\wss{You do not need to report every interaction, but you should disclose the 
uses that materially contributed to your work}

\subsection{Verification of AI-Generated Content}

\wss{\begin{itemize}
    \item Describe how the team checked AI-generated or AI-assisted material
    before including it in the document.
    \item In particular, verify factual claims, technical assertions,
    requirements, calculations, references, and other information for which an
    AI system may produce incorrect or fabricated results.
    \item \textbf{The team is responsible for the accuracy, completeness,
consistency, and appropriateness of the submitted document, regardless of
whether material was generated by a human or an AI system.}
\end{itemize}}

\newpage{}

\section{Appendix --- Reflection}

\wss{Not required for CAS 741 projects}

The information in this section will be used to evaluate the team members on the
graduate attribute of Problem Analysis and Design.

\input{../../Reflection.text}

\begin{enumerate}
  \item What went well while writing this deliverable? 
  \item What pain points did you experience during this deliverable, and how
    did you resolve them?
  \item Which of your design decisions stemmed from speaking to your client(s)
  or a proxy (e.g. your peers, stakeholders, potential users)? For those that
  were not, why, and where did they come from?
  \item While creating the design doc, what parts of your other documents (e.g.
  requirements, hazard analysis, etc), it any, needed to be changed, and why?
  \item What are the limitations of your solution?  Put another way, given
  unlimited resources, what could you do to make the project better? (LO\_ProbSolutions)
  \item Give a brief overview of other design solutions you considered.  What
  are the benefits and tradeoffs of those other designs compared with the chosen
  design?  From all the potential options, why did you select the documented design?
  (LO\_Explores)
  \item (After you have implemented another team's module, which means this
  isn't filled in until after the original deadline). What did you learn by
  implementing another team's module? Were all the details you needed in the
  documentation, or did you need to make assumptions, or ask the other team
  questions?  If your team also had another team implement one of your modules,
  what was this experience like?  Are there things in your documentation you
  could have changed to make the process go more smoothly for when an
  ``outsider'' completes some of the implementation?

  Your reflection should be about concepts learned in your SE education, such as modular decomposition, cohesion, coupling, information hiding, etc.  You should critique the quality and utility of the specification you received.

  You should create a PR for your implemented module.  A link to the PR should be part of your reflection.

\end{enumerate}

\end{document}